\chapter{Analisis de antecedentes y aportacion realizada}\label{analanteced}

\section{Contexto del sector}

El mercado de distribucion alimentaria en Espana presenta una combinacion de gran
competencia entre cadenas, alta sensibilidad al precio y fuerte fragmentacion de oferta.
Durante los ultimos anos, la inflacion en alimentacion ha incrementado la necesidad de
herramientas que permitan comprar mejor sin aumentar el tiempo invertido.

Aunque existen comparadores de precios consolidados, la mayor parte se centra en mostrar
informacion de coste por tienda, sin integrar de forma nativa el impacto logistico de
desplazarse a varios establecimientos.

\section{Estado del arte en comparacion de precios}

La evolucion observada puede resumirse en tres generaciones:

\begin{enumerate}
	\item \textbf{Directorios de precios:} listados estaticos o semiactualizados, con
		poca cobertura y baja frecuencia de refresco.
	\item \textbf{Comparadores multisupermercado:} soluciones como Soysuper, Findit u OCU
		Market, con mejor catalogo y funcionalidades de lista.
	\item \textbf{Personalizacion con IA:} tendencia emergente en la que se introducen
		recomendaciones contextuales y asistencia conversacional.
\end{enumerate}

En Espana, el salto completo hacia la tercera generacion todavia es limitado,
especialmente en el segmento de cesta alimentaria diaria.

\section{Tecnologias relevantes para BargAIn}

\subsection{Geoespacial y rutas}

PostGIS permite resolver consultas de proximidad con precision y rendimiento, mientras que
OR-Tools facilita modelar el problema de ordenacion de paradas como una variante del TSP/VRP
en instancias pequenas (maximo 3-4 paradas en nuestro caso).

\subsection{Ingesta de datos de precio}

La combinacion Scrapy + Playwright es adecuada para catalogos dinamicos de supermercados.
Sin embargo, por robustez de producto, BargAIn no depende de una unica fuente y combina:

\begin{itemize}
	\item scraping programado,
	\item aportaciones de crowdsourcing,
	\item actualizacion manual de comercios verificados.
\end{itemize}

\subsection{OCR y asistente}

Para OCR se adopta Google Cloud Vision API por mejor rendimiento practico en tickets
y fotos reales. Para el asistente se integra Gemini mediante backend proxy, con
guardrails para limitar respuestas al dominio de la compra.

\section{Aportacion diferencial de BargAIn}

La aportacion del proyecto se concreta en cuatro ejes:

\begin{itemize}
	\item \textbf{Optimizacion multicriterio real} (precio + distancia + tiempo) en lugar de
		comparacion aislada de precios.
	\item \textbf{Inclusion de comercio local} mediante portal business y actualizacion de
		precios/promociones.
	\item \textbf{Digitalizacion de lista} por OCR con validacion de usuario.
	\item \textbf{Interaccion conversacional} para consultas complejas de compra.
\end{itemize}

Estas capacidades, integradas en una sola plataforma, posicionan BargAIn en un espacio
de mercado poco cubierto por herramientas actuales.
